\subsection{rMD17 prior weight sensitivity}\label{sec:weight_sweep}

We first vary the prior weight on rMD17 aspirin. This sweep contains 36 runs: three priors, four weights, and three seeds per cell. We evaluate no prior, the Euclidean covariance prior, and the spectral prior at \(\lambda \in \{0.001,0.01,0.1,1.0\}\). At \(H=16\), the no-prior baseline is \(0.417 \pm 0.199\) for this seed subset. The Euclidean prior is best at \(\lambda=0.01\), where it obtains \(0.173 \pm 0.049\); its performance degrades at larger weights, reaching \(0.238 \pm 0.023\) at \(\lambda=0.1\) and \(0.322 \pm 0.101\) at \(\lambda=1.0\). The spectral prior is best at \(\lambda=0.1\), where it obtains \(0.130 \pm 0.080\), and remains similar at \(\lambda=1.0\) with \(0.144 \pm 0.064\). At \(\lambda=0.01\), however, the spectral sweep has a heavy-tailed failure mode, with mean \(1.908 \pm 1.326\).

This sweep supports two conclusions. First, the main rMD17 setting uses the best spectral weight in the sweep, \(\lambda=0.1\), rather than a value that is favorable only by accident. Second, even when each prior is evaluated at its own tuned best, the Euclidean optimum \(0.173\) does not match the spectral optimum \(0.130\). The instability at spectral \(\lambda=0.01\) is a useful caveat: the spectral prior is not monotonic in weight, and small grids can expose heavy-tail behavior. The no-prior mean in this sweep is also higher than the main rMD17 baseline in Table~\ref{tab:rmd17} (\(0.417\) versus \(0.290\)), because the sweep uses only three seeds rather than the 10-seed main experiment. We therefore treat this analysis as evidence about relative weight sensitivity, not as a replacement for the main rMD17 estimate.

\subsection{rMD17 Laplacian construction mode}\label{sec:rmd17_lapmode}

We next vary how the spectral Laplacian is constructed on rMD17 aspirin. This ablation contains 15 runs: three Laplacian modes and five seeds per mode. The three modes are per\_frame, which recomputes the Laplacian for each transition; fixed\_mean, which uses an average Laplacian; and fixed\_frame0, which uses the graph from a single initial frame. At \(H=16\), per\_frame obtains \(0.111 \pm 0.075\), fixed\_mean obtains \(0.131 \pm 0.112\), and fixed\_frame0 obtains \(0.146 \pm 0.089\). All three are below the no-prior main-run baseline of \(0.290\) from Table~\ref{tab:rmd17}.

The differences among the three Laplacian modes are small relative to their seed-to-seed variability. Per\_frame is numerically best, but its standard deviation overlaps substantially with fixed\_mean and fixed\_frame0, so this ablation does not support a strong ordering among modes on aspirin. The result is consistent with the structure of the molecular setting: the covalent bond graph is essentially invariant across MD transitions, since bonds do not break in these trajectories. Consequently, per-frame, mean, and initial-frame Laplacians are nearly the same object. This ablation therefore clarifies that, for rMD17 aspirin, the benefit of the spectral prior is not sensitive to the precise Laplacian construction mode, provided the mode reflects the stable molecular bond graph.

\subsection{Wolfram CA Laplacian construction and weight}\label{sec:wolfram_mechanism}

Table~\ref{tab:wolfram_mechanism} and Figure~\ref{fig:wolfram_mechanism} report the corresponding mechanism ablation on Wolfram cellular automata. This experiment contains 50 runs: a five-seed no-prior baseline plus spectral priors across three Laplacian modes, three weights, and five seeds per cell. The no-prior baseline obtains \(0.066 \pm 0.043\) at \(H=16\). With per-step Laplacian recomputation, the spectral prior is stable at small weights: \(0.098 \pm 0.073\) at \(\lambda=0.001\) and \(0.062 \pm 0.018\) at \(\lambda=0.005\), matching the baseline within seed variability. At \(\lambda=0.01\), per\_step enters a transition regime, with mean \(1.894 \pm 3.522\) driven by one catastrophic seed (8.934); the other four seeds remain in the range \([0.044, 0.348]\). Fixed\_initial is unstable at all weights, increasing from \(4.279 \pm 6.991\) at \(\lambda=0.001\) to \(118.978 \pm 187.912\) at \(\lambda=0.005\) and \(13046.638 \pm 26086.373\) at \(\lambda=0.01\), with a maximum H=16 error of 65219. Fixed\_average is intermediate: it is close to baseline at \(\lambda=0.001\) with \(0.084 \pm 0.056\), but becomes unstable at larger weights.

This ablation refines the interpretation of the Wolfram result. The spectral prior is conditionally stable on Wolfram CA, not fundamentally broken: when the Laplacian is recomputed per step and the weight is small, it matches the no-prior baseline. The failures arise when a static-graph approximation is combined with a stronger prior weight. This differs from rMD17 because the relevant graph in the cellular automaton evolves rapidly across transitions. A fixed\_initial Laplacian can therefore impose a graph smoothness constraint that no longer matches the current state, and increasing \(\lambda\) amplifies this systematic mismatch. Fixed\_average partially reduces this mismatch at the smallest weight but still becomes unstable once the constraint is strengthened. Thus, the Wolfram ablation identifies a mechanism rather than a contradiction: in this setting, spectral regularization is useful when the Laplacian is recomputed to track the evolving cell graph, but static Laplacian approximations destabilize training as the prior weight grows. Whether this pattern generalizes to other domains with non-stationary graph structure remains an open question that we address briefly in Section~\ref{sec:discussion}.
